% Supervisor meeting, 12 Nov 2026 (DS500). Purpose: report the state of dev, explain the
% decisions taken so far (recorded with an ID or implicit in the design principles) with the
% alternative we rejected, get answers to the open questions, and show the outlook up to a
% Germany-wide reach. Short version: 12 main slides, details in the backup. Frozen after the
% meeting; the next meeting gets a new deck (slides/README.md, "Versioning").
% Rules: slides/AGENTS.md. Every number comes from common/facts.tex. Open items are marked as
% open; nothing proposed is shown as approved; the Germany outlook is a vision of the team
% (Juju, 12 Nov 2026), not part of the docs' plan.
\documentclass[aspectratio=169,10pt]{beamer}
\usetheme{reach}
\input{facts}
\input{isochrone}
\input{system}
\input{hosting}

% Marker for decisions without an ID (implicit in the design principles).
\newcommand{\impl}{\textit{implicit}}
% Ragged-right table column (no stretched word spaces in narrow cells).
\newcolumntype{P}[1]{>{\raggedright\arraybackslash}p{#1}}

\title[reach $\cdot$ Supervisor meeting]{reach: our decisions, open questions and the way to
  all of Germany}
\subtitle{Supervisor meeting $\cdot$ DS500 Data Science Project, Topic 9 ``The 15-Minute City''}
\author{Rico $\cdot$ Socke $\cdot$ Juju}
\institute{University of Tübingen, winter term 2026/27}
\date[12 Nov 2026]{12 November 2026}

\begin{document}

{\setbeamertemplate{headline}{}\setbeamertemplate{footline}{}
\begin{frame}
  \titlepage
\end{frame}}

% ============================================================================================
% Key message first (pyramid). Speaker: TBD
\begin{frame}{reach runs end to end for Tübingen; today we explain the decisions behind it, ask
  four questions and sketch the way to all of Germany}
  \centering
  \begin{tikzpicture}[
      box/.style={draw=ink,line width=0.6pt,text width=28.5mm,align=flush left,inner sep=2.2mm,
        font=\footnotesize,minimum height=36mm,anchor=north}]
    \node[fill=ink,text=white,text width=134mm,align=center,inner sep=2.5mm,font=\small\bfseries]
      (key) at (0,0)
      {Freeze the Tübingen scope and deploy it; keep every choice open for a larger area};
    \node[box] (a) at (-51mm,-11mm) {{\usebeamercolor[fg]{alerted text}\bfseries Done}\\[1mm]
      Pipeline, API and web app on \texttt{dev}; \factCurCategories{} categories, 4 profiles;
      recommender, combined score and housing potential; 9 of 11 validation experiments};
    \node[box] (b) at (-17mm,-11mm) {{\usebeamercolor[fg]{alerted text}\bfseries Decided}\\[1mm]
      Open data only, the official 100\,m census grid instead of H3, Valhalla only, counting with
      the matrix, the paper's classes, a greedy recommender};
    \node[box] (c) at (17mm,-11mm) {{\usebeamercolor[fg]{alerted text}\bfseries Open}\\[1mm]
      Hosting (Pages split, ProDesk or VPS), Moran's I as the ``adequate method'', which known
      limits are acceptable};
    \node[box] (d) at (51mm,-11mm) {{\usebeamercolor[fg]{alerted text}\bfseries Next}\\[1mm]
      Submission on 22 January for Tübingen; long-term goal: the same method for all of
      Germany (vision, not planned)};
    \foreach \n in {a,b,c,d} \draw[ink,line width=0.6pt] (key.south) -- (\n.north);
  \end{tikzpicture}
  \source{run \factCurRunId{} (OSM \factCurOsmDate); validation report (run \factValRunId);
    docs/architecture.md §13; docs/deployment\_proposal.md (proposal, PR \#64).}
\end{frame}

% ============================================================================================
\section{Status}
% Speaker: TBD

\begin{frame}{dev now measures \factCurCategories{} categories for every resident of Tübingen
  and publishes them in one validated run}
  \begin{columns}[T,onlytextwidth]
    \column{0.27\textwidth}
      \keyfigure{\factCurCells}{populated 100\,m cells (\factCurResidents{} residents)}
      \par\bigskip
      \keyfigure{\factCurCategories{} $\times$ 4}{categories $\times$ profiles}
    \column{0.27\textwidth}
      \keyfigure{\factCurFacilities}{facilities in the routing area}\par\bigskip
      \keyfigure{\factCurRecs}{recommended zones; \factCurHousing{} housing cells}
    \column{0.42\textwidth}
      \textbf{Merged into \texttt{dev} since 10 November}
      \begin{itemize}\footnotesize
        \item \#57: parks and playgrounds separate; private ones no longer count
        \item \#58: combined score 0--100 over chosen categories (display only)
        \item \#59: housing potential for cells without residents
        \item \#60: eight extended categories in six groups
        \item \#64: hosting proposal (document only, not approved)
      \end{itemize}
      {\footnotesize All \factCurChecks{} result checks pass; three warn.}
  \end{columns}
  \source{run \factCurRunId{} (OSM \factCurOsmDate), validate\_results report part 2; git log of
    \texttt{dev}; docs/features.md.}
\end{frame}

% ============================================================================================
\section{Decisions}
% Speaker: TBD

\begin{frame}{Every layer of reach rests on a deliberate choice; for each we can name the
  alternative we rejected}
  \footnotesize
  \begin{tabular}{@{}P{22mm}P{60mm}P{50mm}@{}}
    \toprule
    Layer & Our choice & Instead of \\
    \midrule
    Data & Open data, one snapshot per run & Public APIs at runtime (Overpass, Nominatim) \\
    Spatial unit & Official Zensus 2022 100\,m grid, cell centroid & H3 hexagons, building
      origins \\
    Study area & Every resident of the city, 4\,km routing buffer & Built-up area only \\
    Routing & Valhalla, self-hosted, one tile set & OSMnx/NetworkX, pandana, r5py \\
    Counting & Time matrix; isochrones only drawn & Point-in-polygon with isochrones \\
    Profiles & Four profiles, 3.39\,km/h default, slopes on & One average speed, flat ground \\
    Scoring & $k_c$, three classes, deficit (Mouratidis 2026) & One stored composite index \\
    Recommender & Greedy on the deficit, 100\,m zones & Exact solver in every run, exact
      sites \\
    System, team & One command, validated publish, docs first & Manual steps, notebooks \\
    \bottomrule
  \end{tabular}
  \par\medskip
  {\scriptsize\color{muted}Recorded decisions carry an ID (O: architecture, M: methodology, W:
  app iteration, N: notes); \impl{} ones follow from a design principle.}
  \source{docs/architecture.md §2, §12, §13; methodology.md; notes.md; AGENTS.md (architecture
    rules).}
\end{frame}

\begin{frame}{We keep the official Zensus 100\,m grid instead of H3 hexagons, so every resident
  count stays exactly as published}
  \footnotesize
  \begin{tabular}{@{}P{20mm}P{57mm}P{55mm}@{}}
    \toprule
    & \textbf{Zensus 100\,m grid (chosen)} & H3 hexagons (rejected) \\
    \midrule
    Residents & Exact counts as published & Re-gridding spreads residents by area; errors at
      every edge \\
    Identity & Official \texttt{GITTER\_ID} (EPSG:3035) joins other census grid tables (age:
      M6) & Own index; every join needs re-gridding \\
    Scale & Same grid for all of Germany: \factDeCells{} populated cells in the file we already
      download & Global, but residents must be disaggregated first \\
    Trade-off & Squares: 8 neighbours at two distances & Hexagons: 6 equal neighbours, nicer
      for spatial statistics \\
    \bottomrule
  \end{tabular}
  \par\medskip
  \textbf{Same grid everywhere:} scored cells, recommender candidates, housing cells and the
  pin-drop snap to it. Origin = cell centroid (O5); with building origins only
  \factVSixChange{} of residents $\times$ categories change class (V6, walking).
  \source{docs/architecture.md §2 (principle 4), §5.2, §12 (H3 aggregation out of scope), O5;
    validation report V6 (run \factValRunId); Zensus file of run \factCurRunId.}
\end{frame}

\begin{frame}{We count along the network with Valhalla's time matrix; isochrone polygons would
  add \factVOneOvercount{} unreachable pairs}
  \begin{columns}[T,onlytextwidth]
    \column{0.55\textwidth}
      \centering\isochronetrap
    \column{0.42\textwidth}
      \footnotesize
      \textbf{\impl:} Valhalla is the only routing engine, self-hosted, one tile set for all
      profiles; map and pin-drop give the same numbers.\par\medskip
      \textbf{D2:} per block of cells one matrix to all facilities in reach; a facility counts
      if its time is $\le 15$ min. Isochrones are only drawn.\par\medskip
      \textbf{V1} (150 cells, walking 3.39\,km/h): polygons add \factVOneOvercount{}
      unreachable pairs, \factVOneBarrier{} of them across the Neckar, rail lines or main
      roads.\par\medskip
      \textbf{O20:} exact prefilter; \factCurFalseNeg{} sampled pairs outside the counts are
      reachable by \texttt{/route}.
  \end{columns}
  \source{AGENTS.md (architecture rule 1); docs/architecture.md §6.3 (D2), §6.5, O3, O20;
    validation report V1 (run \factValRunId); run \factCurRunId{} (\texttt{false\_negatives}).}
\end{frame}

\begin{frame}{We measure inclusively, follow the paper's classes and publish only results that
  pass every check}
  \footnotesize
  \begin{tabular}{@{}P{11mm}P{50mm}P{27mm}P{42mm}@{}}
    \toprule
    ID & Decision & Instead of & Why \\
    \midrule
    O11 & Score every populated cell in the city; the built-up area only shapes map and
      candidates & Built-up area only & That left \factUrbanUnscored{} of residents unscored
      (N3) \\
    O6, O7 & Walking 3.39\,km/h (inclusive) and 5.0, cycling 15, wheelchair; slope model on &
      One speed, flat ground & Tübingen is hilly: slopes change \factVElevenChange{} of
      classes (V11) \\
    O8 & Seven Core categories as in the paper, eight extended ones as indicative & Fine
      categories (GPs, school types) & OSM gaps stay small \\
    §4 & $k_c$, three classes, deficit weighted by residents; combined score display only &
      One stored index & A planner sees what is missing \\
    §5 & Greedy recommender on the deficit, 100\,m zones & Exact solver in every run & At most
      \factVFourGap{} below the optimum (V4) \\
    O19 & One command; \texttt{publish} only after all checks; no check is loosened & Manual
      steps & A failed run never reaches the app (N16) \\
    \bottomrule
  \end{tabular}
  \source{docs/architecture.md §4, §6.1, §7, §13; methodology.md §2, §4, §4b, §5; notes N3, N16;
    validation report V4, V11 (run \factValRunId).}
\end{frame}

% ============================================================================================
\section{Open questions}
% Speaker: TBD

\begin{frame}{Four questions need your answer; for each we bring a recommendation}
  \footnotesize
  \begin{tabular}{@{}P{6mm}P{58mm}P{46mm}P{20mm}@{}}
    \toprule
    \# & Question & Our recommendation & Blocks \\
    \midrule
    1 & Frontend on Cloudflare Pages, backend on one host (proposal of \#64, not approved)? &
      Decide after question 2 & Deployment \\
    2 & Is the HP ProDesk acceptable as backend host, or is a rented VPS required? & Measure the
      first full refresh on the ProDesk, then decide & Server, budget \\
    3 & Moran's I shows strong clusters on cells (\factVEightMoran), none between districts:
      is this the ``adequate method'' (M5)? & Yes: optimization plus spatial statistics; LISA
      map in the app if you wish & Scope of V8 \\
    4 & Which known limits do you accept (wheelchair on slopes, OSM gaps, extended categories
      unchecked)? & State them openly; map missing Kitas and schools in OSM & Final report \\
    \bottomrule
  \end{tabular}
  \par\medskip
  {\footnotesize\textbf{Also:} interim format; database access for you (SSH tunnel?); may we use
  bwUniCluster; is one city enough, and is a larger area of interest?}
  \source{docs/deployment\_proposal.md (PR \#64); methodology.md §7 (M5); validation report V8
    (run \factValRunId); notes N11, N13, N16; details in the backup.}
\end{frame}

% ============================================================================================
\section{Outlook}
% Speaker: TBD

\begin{frame}{Until 22 January we finish Tübingen: close the validations, deploy, and deliver}
  \footnotesize
  \begin{tabular}{@{}P{26mm}P{106mm}@{}}
    \toprule
    When & What \\
    \midrule
    25 Nov & Interim presentation (\texttt{v0.1.0}); demo local unless decided otherwise (O16) \\
    from $\sim$1 Dec & Server live after the hosting decision; weekly refresh by timer \\
    Dec -- Jan & V7 (local knowledge) and V9 (rerun from the README); M1 final categories, M2
      transit classes; recommender decision cards and equity mode (M4) \\
    10 -- 15 Jan & Feature freeze: Stretch only if Core passes validation \\
    22 Jan & Submission (\texttt{v1.0.0}): code, hosted app, validation report, video, summary \\
    29 Jan & Final presentation with live demo \\
    \bottomrule
  \end{tabular}
  \source{docs/features.md §9 (roadmap); project\_management\_decisions.md §8; notes N12, N13;
    methodology.md §7.}
\end{frame}

\begin{frame}{Our long-term goal is all of Germany, and the core choices already scale}
  \footnotesize
  \begin{tabular}{@{}P{26mm}P{50mm}P{56mm}@{}}
    \toprule
    Building block & Today: Tübingen & Germany \\
    \midrule
    Residents & Zensus 100\,m grid, filtered from the nationwide file & Same file:
      \factDeCells{} populated cells, \factDeResidents{} residents \\
    Network, facilities & Geofabrik extract of the Regierungsbezirk Tübingen, clipped & Geofabrik
      extract of Germany, same format and code \\
    Study area & One municipality key (AGS) in the config & All municipalities, one run each,
      in parallel \\
    Routing & Valhalla tiles built on our server & Valhalla routes up to planet scale; build on
      bwUniCluster (§11.1) \\
    Transit & NVBW GTFS (naldo) & A nationwide GTFS feed (to check, e.g.\ DELFI) \\
    \bottomrule
  \end{tabular}
  \par\medskip
  {\footnotesize\color{muted}Vision of the team, not part of the course plan: a second region is
  out of scope for Core (architecture §12). Nothing on this or the next slide is built.}
  \source{Zensus file of run \factCurRunId{} (grid sum, perturbed); config/config.yaml
    (\texttt{study\_area}, \texttt{sources}); docs/architecture.md §4, §11.1, §12.}
\end{frame}

\begin{frame}{Three steps lead to Germany; compute, data quality, comparability and serving
  decide whether it is feasible}
  \centering
  \begin{tikzpicture}[
      stair/.style={draw=ink,line width=0.6pt,text width=38mm,align=flush left,inner sep=1.8mm,
        font=\footnotesize,anchor=south west},
      head/.style={font=\footnotesize\bfseries,text=accent,anchor=south west,inner sep=0.5mm},
      prob/.style={draw=muted,line width=0.5pt,text width=28.5mm,align=flush left,
        inner sep=1.6mm,font=\scriptsize,anchor=north west,minimum height=17mm}]
    \node[stair,minimum height=11mm] (s1) at (0,0) {Every municipality of the extract we
      already download};
    \node[head] at (s1.north west) {1 $\cdot$ Regierungsbezirk Tübingen};
    \node[stair,minimum height=17mm] (s2) at (44mm,0) {State extract; one job per district on
      bwUniCluster; counts cached across runs};
    \node[head] at (s2.north west) {2 $\cdot$ Baden-Württemberg};
    \node[stair,minimum height=23mm] (s3) at (88mm,0) {Nationwide extract, same Zensus file;
      vector tiles and the 1\,km grid for the zoomed-out map};
    \node[head] at (s3.north west) {3 $\cdot$ Germany};
    \node[prob] at (0,-3mm) {{\bfseries Compute.} \texttt{count}: \factCurCountRuntime{} for
      \factCurCells{} cells; \factDeCellFactor{}$\times$ the cells $\approx$
      \factDeCountDays{} days naively};
    \node[prob] at (32.5mm,-3mm) {{\bfseries Data quality.} OSM completeness differs by region;
      V2 needs official lists per region};
    \node[prob] at (65mm,-3mm) {{\bfseries Comparability.} 15 min and $k_c$ fit a city, not a
      village: density strata (M3)};
    \node[prob] at (97.5mm,-3mm) {{\bfseries Serving.} \factCellsPayload{} GeoJSON per layer
      today; vector tiles and more than 20\,\texteuro};
  \end{tikzpicture}
  \source{notes N16 (count runtime, run \factCurRunId); naive linear estimate from the Zensus
    file; implementation\_plan.md P6 (payload); architecture.md §5.2, §9, §11.1; methodology.md
    M3; proposal of the team, not in the docs' plan.}
\end{frame}

{\setbeamertemplate{headline}{}
\begin{frame}
  \vfill
  {\Large\bfseries Thank you. Your decisions, please.\par}
  \vskip4mm
  {\color{accent}\rule{25mm}{2pt}}\par\vskip4mm
  Hosting, the adequate method and the acceptable limits decide what we freeze next;\\
  all of Germany stays our long-term goal.\par
  \vfill
\end{frame}}

% ============================================================================================
\backupslides
\section*{Backup}

\begin{frame}{Backup: question 1, current hosting design vs.\ the proposal of \#64}
  \begin{columns}[T,onlytextwidth]
    \column{0.5\textwidth}
      \centering{\footnotesize\bfseries Current design (binding)}\par\smallskip
      \hostingcurrent
    \column{0.5\textwidth}
      \centering{\footnotesize\bfseries Proposal (not approved)}\par\smallskip
      \hostingproposed
  \end{columns}
  \source{docs/architecture.md §3, §11 (current); docs/deployment\_proposal.md (proposal, PR \#64).}
\end{frame}

\begin{frame}{Backup: question 2, the ProDesk saves rent, a VPS saves operating effort}
  \footnotesize
  \begin{tabular}{@{}P{26mm}P{50mm}P{50mm}@{}}
    \toprule
    & HP ProDesk (i5-8500T, 16\,GB RAM) & Rented VPS ($\ge 8$\,GB, O1) \\
    \midrule
    Benefit & Hardware is available; no rent & Data-centre power and network; remote
      management \\
    To check & Power and network uptime, remote recovery, storage, backups, who maintains it
      & Monthly cost within the 20\,\texteuro{} budget, disk size, provider limits, backups \\
    Evidence needed & First full tile build and refresh measured on the host & Same measurement on
      the chosen plan \\
    \bottomrule
  \end{tabular}
  \par\medskip
  \textbf{Please clarify:} does ``keep VPS and hosting separate'' require a rented VPS? What
  uptime does the course demonstration need? Is a private network for team access (for example
  Tailscale) acceptable, and how should you reach the read-only database account?
  \source{docs/deployment\_proposal.md (proposal, PR \#64); docs/architecture.md §11 (resource
    estimate to be verified on first deployment).}
\end{frame}

\begin{frame}{Backup: question 4, five known limits and what we propose}
  \footnotesize
  \begin{tabular}{@{}P{32mm}P{62mm}P{34mm}@{}}
    \toprule
    Limit & Effect & Proposal \\
    \midrule
    Wheelchair on slopes & Steep sections are not avoided; hillside counts are optimistic (N11)
      & Accept, state on data page \\
    OSM gaps & \factVTwoKitaMissing{} Kitas and \factVTwoSchoolMissing{} schools have no OSM
      match (V2) & Map them in OSM \\
    Extended categories & No completeness check; monotonicity check at \factCurMonoCells{} cells,
      fails above 1\,\% (N16) & Mark as indicative \\
    Cycling recommender off & \factCyclingZones{} zones would be added; Valhalla needs
      \factCyclingValhallaPeak{} (W6) & Accept, document \\
    Housing candidates & Land use not fully excluded; buildable not known (N15) & Decide after
      map check \\
    \bottomrule
  \end{tabular}
  \par\medskip
  {\footnotesize\color{muted}Shares depend strongly on speed and budget (V5): we always name the
  profile and budget next to a share.}
  \source{notes N11, N12, N13, N15, N16; validation report V2, V5 (run \factValRunId); run
    \factCurRunId.}
\end{frame}

\begin{frame}{Backup: five of the eight extended categories are beyond a 15-minute walk for
  nearly half of the residents or more}
  \centering
  \begin{tikzpicture}
    \begin{axis}[
        reach ranking, reach box legend, xbar stacked, bar width=2.4mm, height=58mm,
        yticklabels from table={\factCurClassesWalk}{category}, xmin=0, xmax=100,
        xlabel={Share of residents [\%], walking 3.39\,km/h, 15 min},
      ]
      \addplot[fill=lack,draw=none,
        nodes near coords={\pgfmathfloattofixed{\pgfplotspointmeta}\ifdim\pgfmathresult pt<5pt\else
          \pgfmathprintnumber[fixed,precision=0]{\pgfplotspointmeta}\fi},
        nodes near coords style={font=\tiny\bfseries,text=white,anchor=east,xshift=-0.3mm},
        point meta=explicit]
        table[x=lack,y expr=\coordindex,meta=lack] {\factCurClassesWalk};
      \addplot[fill=insufficient,draw=none]
        table[x=insufficient,y expr=\coordindex] {\factCurClassesWalk};
      \addplot[fill=sufficient,draw=none]
        table[x=sufficient,y expr=\coordindex] {\factCurClassesWalk};
      \legend{Lack, Insufficient, Sufficient}
    \end{axis}
  \end{tikzpicture}
  \source{run \factCurRunId{} (OSM \factCurOsmDate), \factCurResidents{} residents in
    \factCurCells{} cells. ext.\ = extended category, indicative (no completeness check).}
\end{frame}

\begin{frame}{Backup: validation experiments V1--V11}
  \scriptsize
  \begin{tabular}{@{}llP{95mm}@{}}
    \toprule
    \# & Experiment & Key result \\
    \midrule
    V1 & Matrix vs.\ isochrone polygons & Polygons add 18--20\,\% false pairs, 47\,\% across barriers \\
    V2 & OSM completeness & Childcare recall 88\,\%, precision 98\,\%; schools 95\,\% / 95\,\% \\
    V3 & Recommender back-test & 80--98\,\% of picks within 500\,m of a removed facility (only the removed deficit counts) \\
    V4 & Greedy vs.\ exact MILP & At most 2.8\,\% below the optimum (mean 0.17\,\%) \\
    V5 & Sensitivity ($k$, speed, budget) & Shares move strongly; district rankings stable
      (median Spearman 0.78--0.95) \\
    V6 & Centroid vs.\ building origin & 1.9\,\% (walking) / 2.8\,\% (wheelchair) change class \\
    V7 & Local knowledge & Not yet run \\
    V8 & Spatial statistics & Moran's I 0.72--0.88 on cells; none between districts \\
    V9 & Reproducibility & Not yet run \\
    V10 & Matrix vs.\ single routes & Median $|\Delta|$ 0.5\,s; 0 of 900 false negatives \\
    V11 & With vs.\ without slopes & 3.3--3.6\,\% change class, mostly on hillsides \\
    \bottomrule
  \end{tabular}
  \source{docs/validation\_report.md, run \factValRunId{} (Core categories of that run, parks
    and playgrounds still one category); the extended categories are not validated yet.}
\end{frame}

\end{document}
